\documentclass[11pt,a4paper]{article}

%------------------------------------------------
% Packages
%------------------------------------------------

\usepackage[margin=2.5cm]{geometry}
\usepackage[T1]{fontenc}
\usepackage[utf8]{inputenc}
\usepackage{lmodern}

\usepackage{amsmath,amssymb}
\usepackage{booktabs}
\usepackage{longtable}
\usepackage{array}

\usepackage{xcolor}
\usepackage{listings}

\usepackage{graphicx}
\usepackage{subcaption}
\usepackage{float}

\usepackage[hidelinks]{hyperref}
\usepackage{cleveref}

\usepackage{enumitem}
\usepackage{tcolorbox}

\setcounter{tocdepth}{3}

%------------------------------------------------
% Listings style  (Dynare .mod highlighting)
%------------------------------------------------

\lstdefinestyle{dynare}{
    basicstyle=\ttfamily\small,
    keywordstyle=\color{blue},
    commentstyle=\color{green!50!black},
    stringstyle=\color{red!70!black},
    frame=single,
    breaklines=true,
    columns=fullflexible,
    morekeywords={var, varexo, parameters, model, end, shocks, stderr,
                  stoch_simul, estimation, estimated_params, steady, check,
                  varobs, initval, histval, calib_smoother, set_dynare_seed}
}

\lstdefinestyle{matlab}{
    basicstyle=\ttfamily\small,
    keywordstyle=\color{blue},
    commentstyle=\color{green!50!black},
    stringstyle=\color{red!70!black},
    frame=single,
    breaklines=true,
    columns=fullflexible,
    language=Matlab
}

%------------------------------------------------
% Custom environments and commands
%------------------------------------------------

\newtcolorbox{cmdsyntax}{
    colback=blue!3, colframe=blue!40!black,
    fonttitle=\bfseries\ttfamily, boxrule=0.5pt,
    left=6pt, right=6pt, top=4pt, bottom=4pt
}

\newcommand{\optentry}[4]{%
    \paragraph{\texttt{#1}\hfill\textnormal{\textit{#2}}}%
    #4\par\smallskip\noindent Default:~\texttt{#3}.\par\medskip
}

\newcommand{\outputvar}[2]{%
    \paragraph{\texttt{#1}}%
    #2\par\medskip
}

\newcommand{\command}[1]{\texttt{#1}}
\newcommand{\option}[1]{\texttt{#1}}
\newcommand{\field}[1]{\texttt{#1}}
\newcommand{\file}[1]{\texttt{#1}}
\newcommand{\matlab}{\textsc{Matlab}}

%------------------------------------------------
% Title
%------------------------------------------------

\title{\textbf{Learning Expectations in Dynare}\\[0.3em]
       \large Social Learning and Adaptive Learning}
\date{}

%================================================
\begin{document}

\maketitle

\begin{center}
\begin{minipage}{0.85\textwidth}
\small
\textbf{Status.} Working document: on how to use the learning toolbox.
Written by Mathieu Tassel during a research internship at CEPREMAP (2026), as
part of the project adding the \file{+learning} package. 

\medskip
\textbf{References.}
Social learning: Grimaud, Salle \& Vermandel, ``A Dynare Toolbox for Social
Learning Expectations,'' \textit{Journal of Economic Dynamics and Control},
2025.
Adaptive learning: Slobodyan \& Wouters, ``Learning in a Medium-Scale DSGE Model
with Expectations Based on Small Forecasting Models,'' \textit{American Economic
Journal: Macroeconomics}, 4(2), 2012.
\end{minipage}
\end{center}

\tableofcontents
\newpage


%================================================================
%  SECTION 0 — OVERVIEW
%================================================================
\section{Overview}
\label{sec:overview}

The \file{matlab/+learning} package replaces the rational-expectations (RE)
assumption by a mechanism in which agents form expectations from something other
than the model solution. Two mechanisms are available:

\begin{itemize}[leftmargin=2em]
\item \textbf{Social learning} (\file{+learning/+social}). A population of
      agents holds heterogeneous beliefs about the forward-looking variables.
      Each period a fraction of them receives news, beliefs are ranked by a
      discounted forecast-error criterion, and pairs of agents adopt the fitter
      of the two. The aggregate belief wedge enters the policy function through
      an additional matrix $R$. Both simulation and estimation are supported.
\item \textbf{Adaptive learning} (\file{+learning/+adaptive}). Agents forecast
      the forward variables with a small perceived law of motion whose
      coefficients are updated by constant-gain recursive least squares. The
      actual law of motion is re-solved every period from those beliefs.
      Simulation only.
\end{itemize}

Both mechanisms are configured through naming conventions on parameters
declared in the \file{.mod} file.\\

You can find three  examples in \file{examples/learning/}: a social learning
simulation, a social learning estimation, and an adaptive learning simulation.

\medskip
\noindent Sections
\ref{sec:declaration}--\ref{sec:estimation} describe social learning, which is
the more complete of the two. \Cref{sec:adaptive} describes adaptive learning.
\Cref{sec:limitations} lists what is not supported, and
\cref{app:architecture} documents the internal architecture.


%================================================================
%  SECTION 1 — MODEL DECLARATION
%================================================================
\section{Model declaration}
\label{sec:declaration}

Declaring a social learning model requires three things in the \file{.mod}
file:

\begin{enumerate}
    \item \texttt{SL\_active = 1}, so that Dynare routes to the learning path.
    \item One \texttt{SL\_learn\_<var> = 1} for each learned forward variable.
    \item For each of them: \texttt{SL\_mutation\_std\_<var>},
          \texttt{SL\_fitness\_discount\_<var>}, \texttt{SL\_mutation\_prob\_<var>},
          and a news shock \texttt{a\_<var>} in \command{varexo}. None of these
          has a default value. The global options, such as
          \texttt{SL\_n\_agents}, do have one.
\end{enumerate}

%%%%%%%%%%%%%%%%%%%%%%%%%%%%%%%%%%%%%%%%%%%%%%%%%%
\subsection{\texttt{SL\_}-prefixed parameters}
\label{sec:declaration:params}

\subsubsection{Naming convention}

\begin{longtable}{>{\ttfamily\small}p{5.4cm} >{\ttfamily\small}p{3.4cm} p{1.6cm} p{3.6cm}}
\toprule
\normalfont\textbf{Parameter in \file{.mod}}
    & \normalfont\textbf{Field in \texttt{options\_.\allowbreak learning.social}}
    & \normalfont\textbf{Default}
    & \normalfont\textbf{Description} \\
\midrule
\endhead
SL\_active                & ---        & 0     & Master switch. Must be set to 1 explicitly to use learning\\
SL\_learn\_VAR            & id\_SL     & ---   & 1 marks \texttt{VAR} as learned, 0 keeps it at RE. A forward variable with no flag is not learned \\
SL\_n\_agents             & N\_agents  & 300   & Population size $J$\\
SL\_wedding\_proportion   & wed\_p     & 1.0   & Share of the population entering the tournament \\
SL\_t\_max                & t\_max     & 1     & Horizon of the reported $h$-step-ahead expectations \\
\midrule
SL\_mutation\_prob\_VAR       & mut\_p     & ---  & News probability $\mu$ for \texttt{VAR}, must be in $(0,1]$ \\
SL\_mutation\_std\_VAR    & mut\_sd(iv)   & ---  & Belief dispersion $\xi_v$ for \texttt{VAR}, must be positive \\
SL\_fitness\_discount\_VAR & rhoGN(iv)    & ---  & Fitness memory $\delta_v$, must be in $(0,1)$ \\

\bottomrule
\end{longtable}

\medskip
\noindent
The triplet (\texttt{mutation\_prob}, \texttt{mutation\_std},
\texttt{fitness\_discount}) is attached to a forward-looking endogenous variable
by the suffix \texttt{<VAR>}. Declaring a variable as learned without providing
these three values raises an error, and so does a value outside its admissible
range.\\
On the other hand a per-variable parameter attached to a variable that is not
learned has no effect: it raises \texttt{learning:social:orphanParam} as a
warning, which in estimation avoids a flat direction in the likelihood.

\medskip
\noindent
The scan stores the parameters and options in two places:

\begin{itemize}[leftmargin=2em]
\item \field{M\_.params} 
\item \field{options\_.learning.social}
\end{itemize}

The SL hyperparameters can be estimated by Bayesian MCMC like any other
parameter. They are declared in \command{estimated\_params} with a prior, and at
each likelihood evaluation \texttt{refresh\_learning\_params} re-reads
\field{M\_.params} into \field{options\_.learning.social}
(\cref{sec:estimation:likelihood}).

\subsubsection{The \texttt{SL\_active} switch}

\texttt{SL\_active} is the only way to activate social learning. If it is not
declared, or set to 0, Dynare runs on the RE path. The dispatch points in
\file{stoch\_simul.m} and \file{dynare\_estimation\_1.m} read
\texttt{SL\_active} from \field{M\_.params} before calling into the package, so
\texttt{setup\_options\_SL} is never reached with the mechanism off.

\medskip
\noindent
When the mechanism is off, the dispatch calls
\texttt{learning.social.warn\_wedge\_variance}, which reports any
\texttt{a\_VAR} shock still carrying a non-zero variance. Such a run is ``RE
plus an expectation shock'' rather than pure RE.

\subsubsection{Example}

\begin{lstlisting}[style=dynare, caption={SL parameters in the \texttt{.mod} file}]
parameters  beta kappa sigma phi_pi phi_y rho_mp        // structural
            SL_active SL_n_agents                       // SL global
            SL_learn_pi                                 // SL selection
            SL_mutation_std_pi SL_fitness_discount_pi SL_mutation_prob_pi ;  // SL per-variable

beta   = 0.99;  kappa = 0.03;  sigma  = 1;
phi_pi = 1.50;  phi_y = 0.125; rho_mp = 0;

SL_active              = 1;      // required: nothing runs without it
SL_n_agents            = 300;
SL_learn_pi            = 1;      // required: selects the learned set
SL_mutation_prob_pi    = 0.3;
SL_mutation_std_pi     = 0.003;
SL_fitness_discount_pi = 0.5;
\end{lstlisting}


%%%%%%%%%%%%%%%%%%%%%%%%%%%%%%%%%%%%%%%%%%%%%%%%%%
\subsection{Belief shocks}
\label{sec:declaration:beliefs}
 
A learned variable enters the model through an expectation equation of the
form 
 
\begin{lstlisting}[style=dynare]
pie = pi(+1) + a_pi;
\end{lstlisting}
 
\noindent
where \texttt{pie} is the subjective expectation of next-period inflation, the
\texttt{(+1)} term is the model-consistent component, and \texttt{a\_pi} carries
the deviation between the two. \texttt{a\_pi} is determined by the learning
process.

\noindent
With \texttt{var a\_pi = 0}, \texttt{a\_pi} carries only the belief wedge: an
object produced by the learning mechanism, never drawn. It is declared as a
\command{varexo} because that is where the implementation stores it.

\noindent
Giving it a positive variance adds an exogenous news shock on top. Unlike an
ordinary expectation shock, it does not just move expectations for one period,
it enters the beliefs of the agents who receive it, and part of it survives
into later periods through the agents the tournament keeps.
 
\subsubsection{The variance of the news shock}
\label{sec:declaration:newsvariance}

In simulation, the variance is a modelling choice. A zero variance and a
strictly positive one describe two different models:

\begin{itemize}[leftmargin=2em]
\item \emph{Pure social learning}: leave \texttt{a\_VARNAME} out of the
      \command{shocks} block, or set its variance to zero. Beliefs then
      deviate from rational expectations only through what the mechanism
      generates on its own. 
\item \emph{Learning with sentiment}: give it a positive variance. Beliefs
      then have an exogenous driver alongside the endogenous one.
\end{itemize}

In estimation, the variance must be strictly positive, the lower
bound of its prior must also be strictly positive . The likelihood evaluates the Gaussian density of the shocks recovered by inversion; a zero-variance shock makes that density degenerate, and
the mode-finder drives $\sigma_{a_v}$ onto its lower bound.


\subsubsection{The selection matrices \texttt{Or} and \texttt{Ox}}

The scan produces two matrices in \field{options\_.learning.social}:

\begin{itemize}[leftmargin=2em]
\item \field{Or} ($n_{\text{fwrd}} \times n_{\text{exo}}$): entry $(i,j) = 1$
      if the $i$-th forward variable is learned and the $j$-th exogenous is its
      belief channel. 
\item \field{Ox} ($n_{\text{exo}} \times n_{\text{exo}}$): the complementary
      projector $I - O_r'O_r$, which selects the structural shocks. 
\end{itemize}

%================================================================
%  SECTION 2 — STOCHASTIC SIMULATION
%================================================================
\section{Stochastic simulation under Social Learning}
\label{sec:simulation}

When \field{options\_.learning.social} is populated, \command{stoch\_simul}
activates the SL pipeline. The RE decision rules are solved first, then the learning matrix $R$
is computed, then the simulation runs forward with social learning. 

\begin{enumerate}[leftmargin=2em]
\item \texttt{learning.social.setup\_options\_SL} scans the parameters, builds
      \field{options\_.learning.social} and rejects unsupported configurations
      (\cref{sec:simulation:setup_options}).
\item \texttt{learning.social.compute\_ghr} computes $R$ from the dynamic
      Jacobian (\cref{sec:simulation:compute_ghr}).
\item \texttt{learning.social.simult\_SL} runs the forward simulation with
      mutation and tournament (\cref{sec:simulation:simult_SL}).
\item \texttt{learning.social.moments\_SL} and \texttt{learning.social.irf\_SL}
      produce the Monte Carlo moments and the generalised IRFs
      (\cref{sec:simulation:moments,sec:simulation:irfs}).
\end{enumerate}


%%%%%%%%%%%%%%%%%%%%%%%%%%%%%%%%%%%%%%%%%%%%%%%%%%
\subsection{Standard Dynare options}
\label{sec:simulation:std_options}

The SL path reuses the native \command{stoch\_simul} options.

\optentry{order}{INTEGER}{1}{%
    Order of the Taylor approximation. Social learning requires
    \option{order=1}; a higher order raises
    \texttt{learning:social:order}.}

\optentry{periods}{INTEGER}{0}{%
    Length of each simulated path. Required: SL moments are always simulated,
    never analytical. }

\optentry{drop}{INTEGER}{100}{%
    Burn-in discarded before moments are measured, and common burn-in before
    the impulse in the generalised IRFs.}

\optentry{replic}{INTEGER}{1}{%
    Number of Monte Carlo replications over agent draws, for both moments and
    IRFs. This is what the empirical confidence bands are computed from.}

\optentry{irf}{INTEGER}{40}{%
    Horizon of the generalised IRFs.}

\optentry{ar}{INTEGER}{5}{%
    Number of autocorrelation lags reported.}

\optentry{noprint}{BOOLEAN}{false}{%
    Suppress the summary tables.}

%%%%%%%%%%%%%%%%%%%%%%%%%%%%%%%%%%%%%%%%%%%%%%%%%%
\subsection{Option scanning and rejected configurations}
\label{sec:simulation:setup_options}

\texttt{setup\_options\_SL} is called from the dispatch point at the very
beginning of \command{stoch\_simul}, and only when \texttt{SL\_active} is
non-zero.

Beyond building the option structure, it builds the set of learned variables
from the \texttt{SL\_learn\_VAR} flags and collects their parameters. Every
learned variable must have its $\xi_v$, its $\delta_v$, its $\mu_v$ and its
\texttt{a\_VARNAME} news shock. A missing value, or a value outside its
admissible range, raises an error.

It also calls \texttt{learning.social.check\_support}, which rejects the
configurations listed in \cref{sec:limitations}.


%%%%%%%%%%%%%%%%%%%%%%%%%%%%%%%%%%%%%%%%%%%%%%%%%%
\subsection{Random number streams}
\label{sec:simulation:rng}

Social learning uses random streams for two purposes: the belief path (which
agents mutate, how they are paired) and the structural shocks. The two are
drawn from separate, non-overlapping substreams of a single seed, built by
\texttt{learning.learning\_streams}:

\begin{itemize}[leftmargin=2em]
\item \texttt{'sim'} mode, replication $k$: beliefs on substream $2k-1$, shocks
      on substream $2k$. 
\item \texttt{'estim'} mode: beliefs on a fixed substream, so that the belief
      realisation is frozen across parameter draws.
\end{itemize}

Substreams require \texttt{mrg32k3a} or \texttt{mlfg6331\_64}. If you don't declare one, \texttt{setup\_options\_SL} sets \texttt{mrg32k3a} and prints a
warning.


%%%%%%%%%%%%%%%%%%%%%%%%%%%%%%%%%%%%%%%%%%%%%%%%%%
\subsection{Computing the $R$ matrix (\texttt{compute\_ghr})}
\label{sec:simulation:compute_ghr}

Under social learning the policy function is

\begin{equation}
\label{eq:policy}
y_t = \bar{y} + P\,(y_{t-1} - \bar{y})
      + Q\,O_x\,\varepsilon_t
      + R_f\,\bigl(\varphi_t + O_r\,\varepsilon_t\bigr),
\end{equation}

where $P$ and $Q$ are the RE decision-rule matrices (\field{ghx},
\field{ghu} re-indexed in declaration order), $R_f$ collects the
forward-variable columns of $R$, and $\varphi_t$ is the aggregate belief wedge.
The $O_x$/$O_r$ split is what separates the structural shocks from the belief
channel.
\noindent
Solving the model gives

\begin{equation}
R = -\bigl(I + A^{-1}F\bigr)^{-1} A^{-1} F,
\qquad A \equiv F P + G .
\end{equation}

$R$ is $n_{\text{endo}} \times n_{\text{endo}}$ but only the forward-variable
columns are non-zero. It is computed once after \texttt{resol} and stored in
\field{oo\_.dr.ghr}, leaving \field{ghx}, \field{ghu} and \field{ys}
untouched.


%%%%%%%%%%%%%%%%%%%%%%%%%%%%%%%%%%%%%%%%%%%%%%%%%%
\subsection{Forward simulation (\texttt{simult\_SL})}
\label{sec:simulation:simult_SL}

\paragraph{Algorithm.} At each period $t$, in this order:

\begin{enumerate}[leftmargin=2em]
\item \textbf{Mutation.} For each forward variable $v$, draw
      $\mathrm{round}(\mu_v J)$ agents uniformly, with one independent
      \texttt{randperm} per variable, and perturb their belief by
      $\mathcal{N}(0, \xi_v^2)$. This gives the candidate beliefs $m_{j,t}$.
\item \textbf{Aggregation and transition.} The candidate aggregate
      $\varphi_t = \frac{1}{J}\sum_j m_{j,t}$ enters the policy function
      \eqref{eq:policy}, producing $y_t$.
\item \textbf{News injection.} The realisation of the belief shock is added to
      the mutants' beliefs, rescaled by $1/\mu_v$
      (\cref{sec:simulation:conventions}).
\item \textbf{Fitness.} Each agent's belief is scored by a discounted sum of
      squared forecast errors, the forecast being the one implied by the agent's
      own belief and the policy function.
\item \textbf{Tournament.} Agents are paired at random; both members of a pair
      adopt the belief of the fitter one. This gives the post-tournament
      beliefs $a_{j,t}$, which are the starting point of period $t+1$.
\end{enumerate}

Aggregation precedes the fitness evaluation: beliefs influence $y_t$, and are
then scored against that same $y_t$.

\paragraph{Outputs.}
\begin{itemize}[leftmargin=2em]
\item \field{oo\_.endo\_simul}: simulated path,
      $n_{\text{endo}} \times T$ (no initial-condition column).
\item \field{oo\_.learning.social.beliefs}: post-tournament beliefs
      $a_{j,t}$, $n_{\text{fwrd}} \times J \times T$.
\item \field{oo\_.learning.social.beliefs\_candidate}: post-injection
      candidates $m_{j,t}$, same size.
\item \field{oo\_.learning.social.fitness}: fitness values, same size.
\item \field{oo\_.learning.social.expectations}: aggregate wedge entering
      the policy function, $\varphi_t + O_r\varepsilon_t$,
      $n_{\text{fwrd}} \times T$.
\item \field{oo\_.learning.social.expectations\_ahead}: $h$-step-ahead
      aggregate expectations, $n_{\text{fwrd}} \times T \times \mathtt{t\_max}$.
      Diagnostic only.
\item \field{oo\_.learning.social.final\_beliefs}: beliefs at $t = T$,
      $n_{\text{fwrd}} \times J$.
\item \field{oo\_.learning.social.y\_RE}: the RE trajectory under the same
      structural shocks, for comparison.
\end{itemize}



%%%%%%%%%%%%%%%%%%%%%%%%%%%%%%%%%%%%%%%%%%%%%%%%%%
\subsection{Moments}
\label{sec:simulation:moments}

\texttt{learning.social.moments\_SL} runs \option{replic} independent
simulations and averages the
resulting moments. Every replication starts agents at zero beliefs and
discards \option{drop} burn-in periods.

\paragraph{Outputs.}
\begin{itemize}[leftmargin=2em]
\item \field{oo\_.mean}, \field{oo\_.var},
      \field{oo\_.autocorr\{1:ar\}}: Monte Carlo averages, in the standard
      Dynare format but always simulated.
\item \field{oo\_.learning.social.moments.mean\_ci}, \field{.std\_ci},
      \field{.autocorr\_ci\{1:ar\}}: empirical $2.5/97.5$ percentile
      intervals across replications.
\item \field{oo\_.learning.social.moments.cross\_section}: cross-sectional
      belief statistics: \field{std\_mean}, \field{iqr\_mean},
      \field{skew\_mean} averaged over time and replications, plus the
      time-series versions \field{*\_ts}.
\item \field{oo\_.learning.social.moments.mc\_raw}: the raw replication-level
      data.
\end{itemize}

\paragraph{Display.} \texttt{disp\_moments\_SL} prints three tables: aggregate
moments with their intervals against the RE theoretical standard deviations,
first-order autocorrelations against their RE counterparts, and the
cross-sectional belief moments. The RE column comes from the native
\texttt{th\_autocovariances}.


%%%%%%%%%%%%%%%%%%%%%%%%%%%%%%%%%%%%%%%%%%%%%%%%%%
\subsection{Impulse response functions}
\label{sec:simulation:irfs}

\texttt{learning.social.irf\_SL} computes generalised IRFs by paired Monte
Carlo. For each replication, two simulations are run from the same belief
substream, one baseline and one with an impulse at the first post-burn-in
period, the stream state being saved and restored between them. The GIRF is the
difference, and the reported response is the median across replications.

The impulse is the corresponding column of the lower Cholesky factor,
$cs(:,j)$, exactly as in the native \texttt{irf} path, so the responses are
orthogonalised and comparable with the RE benchmark.

\paragraph{Outputs.} For each shock,
\field{oo\_.learning.social.irf.\textit{shock}} holds \field{median},
\field{mean} and \field{bands}. The median GIRF is also written to the
\field{oo\_.irfs} structure, so Dynare's plotting routines work unchanged.


%%%%%%%%%%%%%%%%%%%%%%%%%%%%%%%%%%%%%%%%%%%%%%%%%%
\subsection{Historical decomposition}
\label{sec:simulation:decomp}

With \field{options\_.learning.social.decomp} set, \texttt{simult\_SL} also
propagates one trajectory per channel: one per structural shock and one per
forward variable for the belief channels. The contributions sum to
$y_t - \bar{y}$ by construction.

\texttt{learning.social.plot\_decomp} draws the stacked contributions with the
total trajectory overlaid.

\medskip
\noindent
The adding-up is exact only when $\mathrm{round}(\mu_v J) = \mu_v J$, for the
reason given in \cref{sec:simulation:conventions}; the residual is reported.


%================================================================
%  SECTION 3 — ESTIMATION
%================================================================
\section{Estimation under Social Learning}
\label{sec:estimation}

The SL likelihood plugs into Dynare's native Bayesian estimation
infrastructure. All standard \command{estimation} options, the mode-finders,
the MCMC and the posterior diagnostics are available.


%%%%%%%%%%%%%%%%%%%%%%%%%%%%%%%%%%%%%%%%%%%%%%%%%%
\subsection{Setting up estimation}
\label{sec:estimation:setup}

Data are declared exactly as in a standard estimation. The binding constraint
is the inversion identification condition: \textbf{the number of observables
must equal the number of exogenous shocks}, belief shocks included. This is the
same requirement as Dynare's \option{conditional\_likelihood}, and it also
implies no measurement error and no missing observations.

\subsubsection{Priors}

\begin{lstlisting}[style=dynare, caption={Priors for SL parameters}]
estimated_params;
  // structural
  rho_d,                  , 0, 0.999, beta_pdf,       0.5,  0.2;
  kappa,                  , 0, 0.999, gamma_pdf,      0.3,  0.05;
  phi_pi,                 , 1, 10.00, gamma_pdf,      1.5,  0.1;
  // shock standard deviations, including the belief channel
  stderr e_d,             , 0, 1,     inv_gamma2_pdf, 0.001, 1;
  stderr e_s,             , 0, 1,     inv_gamma2_pdf, 0.001, 1;
  stderr a_pi,            , 0, 1,     inv_gamma2_pdf, 0.001, 1;
  // SL hyperparameters. 
  SL_fitness_discount_pi, , 0, 0.999, beta_pdf,       0.8,  0.1;
end;
\end{lstlisting}

\noindent
Only parameters attached to a learned variable can be estimated. Bound for $\mu$ and
$\delta$ must be away from the edges of $(0,1)$: a draw outside the admissible range is
rejected with an infinite objective rather than crashing, but a prior that puts
mass there wastes evaluations.

\medskip
\noindent
Some SL parameters are harder to identify because the likelihood is not
continuous in them. This is the case for $\xi$: with a finite population, small
changes in $\xi$ flip the outcome of the tournament, so the log-posterior jumps.
A mode-finder then returns the highest point of that noise rather than a
maximum, and the Hessian at the mode is not usable. One solution is to calibrate
$\xi$ and estimate $\delta$.


\subsubsection{\texttt{setup\_learning\_estimation}}

Before the likelihood loop, \texttt{setup\_learning\_estimation} runs once. It
is the SL counterpart of \texttt{dynare\_estimation\_init} and builds the
observation selection matrix $O$, the forward-variable indices, and the index
maps linking positions in \field{M\_.params} to the \field{mut\_sd},
\field{rhoGN} and \field{mut\_p} vectors. These maps let
\texttt{refresh\_learning\_params} inject the current draw at each evaluation
without re-scanning \field{M\_.param\_names}.


%%%%%%%%%%%%%%%%%%%%%%%%%%%%%%%%%%%%%%%%%%%%%%%%%%
\subsection{The inversion filter (\texttt{IF\_SL})}
\label{sec:estimation:filter}

Given a parameter vector $\theta$ and the solved model:

\begin{enumerate}[leftmargin=2em]
\item Initialise $J$ agents at zero beliefs and the state at its steady state.
\item For each period $t$:
      \begin{enumerate}
      \item run the mutation step to obtain the candidate beliefs and the
            aggregate $\varphi_t$;
      \item invert the observation equation to extract the shocks,
            $\hat{\varepsilon}_t = \lambda^{-1}\bigl(y_t^{\text{obs}} - c
            - O(P\hat{y}_{t-1} + R_f\varphi_t)\bigr)$;
      \item feed $\hat{\varepsilon}_t$ into the transition and into the belief
            recursion, then run the fitness and tournament steps.
      \end{enumerate}
\item Accumulate the per-period log-likelihood contributions.
\end{enumerate}

The inversion Jacobian is
$\lambda = O\,Q\,O_x + O\,R_f\,O_r$, constant across periods. It is checked to
be well conditioned: an ill-conditioned $\lambda$ means the observables carry
almost no independent information about the shocks.

\paragraph{Measurement constant.} The data are compared to
$c + O(y_t - \bar{y})$ with $c = O\bar{y}$, or $c = 0$ when
\option{prefilter} is set and the data have been demeaned. The constant that
was removed is reported in \field{oo\_.Smoother.Constant}, following the native
convention.

\paragraph{Random numbers.} \texttt{IF\_SL} does not build its own stream: the
belief stream is passed as the fifth argument. The caller
(\texttt{dsge\_likelihood\_SL}) builds it in \texttt{'estim'} mode at every
evaluation, which freezes the belief realisation across parameter draws.

\paragraph{Early stopping.} If the cumulated inversion residual exceeds a
threshold, the loop stops and reports how many periods were left unprocessed.
The likelihood is then penalised, and the smoother marks the unreached periods
as \texttt{NaN} rather than leaving them at their initial value.


%%%%%%%%%%%%%%%%%%%%%%%%%%%%%%%%%%%%%%%%%%%%%%%%%%
\subsection{Log-likelihood (\texttt{dsge\_likelihood\_SL})}
\label{sec:estimation:likelihood}

\texttt{learning.social.dsge\_likelihood\_SL} has the same argument and output
signature as the native \texttt{dsge\_likelihood}.

\paragraph{Formula.} With $n_e$ shocks of positive variance and $T$
observations after the presample,

\begin{equation}
\log \mathcal{L}(\theta) =
  -\frac{n_e T}{2}\log(2\pi)
  - \frac{T}{2}\log\bigl|\Sigma_\varepsilon\bigr|
  - \frac{1}{2}\sum_{t}\hat{\varepsilon}_t'\Sigma_\varepsilon^{-1}\hat{\varepsilon}_t
  - T \log\bigl|\det \lambda\bigr| .
\end{equation}

The last term is the Jacobian of the change of variables from observables to
shocks.

\paragraph{Parameter refresh.} At each evaluation,
\texttt{refresh\_learning\_params} copies the current \field{M\_.params} into
the SL option fields using the precomputed index maps. The refreshed values are then range-checked, and
a draw outside the admissible set returns an infinite objective rather than
propagating into the belief loop.


%%%%%%%%%%%%%%%%%%%%%%%%%%%%%%%%%%%%%%%%%%%%%%%%%%
\subsection{Smoother output}
\label{sec:estimation:smoother}

\paragraph{\texttt{store\_SL\_smoother}} routes the inversion filter output
into \field{oo\_}:

\begin{itemize}[leftmargin=2em]
\item \field{oo\_.SmoothedShocks.\textit{exo}}: smoothed shocks, belief
      channels included.
\item \field{oo\_.SmoothedVariables.\textit{endo}}: smoothed state
      trajectory.
\item \field{oo\_.Smoother.SteadyState} and \field{oo\_.Smoother.Constant}, for
      reporting compatibility.
\item \field{oo\_.learning.social.smoother}: the SL-specific payload,
      \field{expectations}, the aggregate belief
      wedge, \field{forward\_states}, \field{forward\_vars}, and the per-shock news series.
\end{itemize}

Periods the inversion did not reach are \texttt{NaN} in every output.

\paragraph{\texttt{plot\_SL\_smoother\_results}} produces four figures:
smoothed structural shocks; smoothed observables overlaid on the data;
the aggregate news channel; and the aggregate belief wedge on the learned
variables, which is the deviation of subjective expectations from their
rational counterpart.

%%%%%%%%%%%%%%%%%%%%%%%%%%%%%%%%%%%%%%%%%%%%%%%%%%
\subsection{Conventions that are not in the papers}
\label{sec:simulation:conventions}

A few implementation choices cannot be derived from the published description
and matter when reading the code.

\paragraph{News rescaling by $1/\mu$.} The aggregate news $O_r\varepsilon_t$ is
injected only into the beliefs of the mutating agents, divided by $\mu_v$, so
that the population average moves by exactly $O_r\varepsilon_t$. 


\paragraph{Truncation of the fitness memory.}
The fitness criterion is a discounted sum of past forecast errors with weights
$\delta_v^k$. The implementation truncates it where the weight falls below
$0.01$, with a maximum of 500 lags.

\begin{lstlisting}[style=matlab]
rho_mat = rhoGN .^ (1:500);
rho_mat = fliplr(rho_mat(:, 1:max(sum(rho_mat > 0.01, 2))))';
T0      = size(rho_mat, 1);
\end{lstlisting}

$T_0$ is also the burn-in of the belief recursion: the loop starts at
$T_0 + \texttt{maximum\_lag} + 1$.


%================================================================
%  SECTION 4 — ADAPTIVE LEARNING
%================================================================
\section{Adaptive learning}
\label{sec:adaptive}

The \file{+learning/+adaptive} package implements constant-gain adaptive
learning. It is a simulation proof of concept, meant to show that a new
learning mechanism can be added with only a few functions and hooks in Dynare.
There is no likelihood and no estimation path for now.

\paragraph{Relation to the reference.}
The algebra and the \matlab{} implementation of this package were inspired by
the article and the replication code of Slobodyan and Wouters (2012). In their
framework beliefs are updated with a Kalman filter, and the constant-gain
specification appears only in their sensitivity analysis. It is that
constant-gain variant which is implemented here. Their projection facility,
which prevents belief updates from destabilising the actual law of motion, is
also reflected in \texttt{simult\_AL}: when the re-solved transition matrix has
a spectral radius above the threshold, the whole update is rejected and the
beliefs, the moment matrix and the transition matrices are kept at their last
stable values. No block of code was copied from the original replication files;
the implementation was written independently from the methodology and the
algebra described in the reference.

\paragraph{Mechanism.} Agents forecast the forward variables with a small
perceived law of motion (PLM), either an intercept only or an intercept plus an
AR(1) slope, estimated variable by variable. Every period the coefficients are
updated by constant-gain recursive least squares, and the actual law of motion
is re-solved from the updated beliefs, so the transition matrices
$(\mu_t, T_t, R_t)$ move over time.


\paragraph{Declaration.} No belief shock is needed: beliefs enter by re-solving
the model, not through an exogenous channel. Only \command{parameters} are
declared.

\begin{longtable}{>{\ttfamily\small}p{4.6cm} >{\ttfamily\small}p{3.2cm} p{1.6cm} p{4.6cm}}
\toprule
\normalfont\textbf{Parameter} & \normalfont\textbf{Field} &
\normalfont\textbf{Default} & \normalfont\textbf{Description} \\
\midrule
\endhead
AL\_active        & ---           & 0    & Master switch \\
AL\_gain          & gain          & 0.02 & Constant gain of the RLS updating \\
AL\_pj\_threshold & pj\_threshold & 1.0  & Projection facility threshold on $\max|\mathrm{eig}(T)|$ \\
AL\_plm\_type     & plm\_type     & 1    & 1 = intercept + AR(1), 0 = intercept only \\
AL\_learn\_VAR    & id\_fwrd      & ---  & 1 marks \texttt{VAR} as learned, 0 keeps it at RE. A forward variable with no flag is not learned; declaring none while \texttt{AL\_active} $=1$ is an error \\
\bottomrule
\end{longtable}

\noindent
The selection rule is identical to the social-learning one: \texttt{AL\_active
= 1} and at least one \texttt{AL\_learn\_<var> = 1} must be declared, otherwise
the run stays on the RE path.


\paragraph{Projection facility.} Under adaptive learning the transition matrix
$T_t$ is re-solved every period from the current beliefs, so nothing guarantees
it stays stable. If agents come to perceive a very high persistence, the
resulting law of motion can have an eigenvalue above one, and the simulated
path diverges.

The projection facility prevents that. At each period it compares
$\max|\mathrm{eig}(T_t)|$ to \texttt{AL\_pj\_threshold}; if the threshold is
exceeded, the new law of motion is discarded and the previous, accepted one is
reused for that period. The beliefs themselves are not modified, they keep
updating freely: only the solution is rejected.

Its purpose is practical, since without it a single unlucky draw can make the
model diverge. Setting the threshold very high, or to \texttt{Inf}, disables the
facility.

\paragraph{Outputs.} \field{oo\_.endo\_simul} plus
\field{oo\_.learning.adaptive.info} with \field{beta\_hist}
($n_{\text{fwrd}} \times n_b \times T$, the belief path), \field{beta\_final}
and \field{pf\_hits} (the number of rejected updates). Monte Carlo moments are
produced by \texttt{moments\_AL} in the same format as the SL ones, with the
belief-coefficient dynamics replacing the cross-sectional statistics.

\paragraph{Not implemented.} Impulse responses: \command{stoch\_simul} warns and
falls back to the RE responses, which ignore the learning dynamics. Estimation,
richer PLMs (full state, AR(2)) and Kalman filter learning are not implemented
either.


%================================================================
%  SECTION 5 — LIMITATIONS
%================================================================
\section{Limitations and rejected configurations}
\label{sec:limitations}

\texttt{check\_support} is called on both the simulation and the estimation
path and raises an error for the configurations below, rather than letting them
run to completion and return wrong results.

\begin{longtable}{p{5cm} p{9.5cm}}
\toprule
\textbf{Configuration} & \textbf{Reason} \\
\midrule
\endhead
\option{order} $> 1$ & The expectations matrix is derived from the first-order
solution. \\
\option{loglinear} & The Jacobian is evaluated at \field{oo\_.dr.ys}, which
\option{loglinear} replaces by its logarithm. \\
Deterministic exogenous variables & The Jacobian slicing does not cover them. \\
Measurement errors (estimated or calibrated) & The inversion filter requires as
many shocks as observables, with no measurement noise. \\
Estimated shock correlations & The $O_x$/$O_r$ split assumes news and
structural shocks are orthogonal. \\
Skewed shocks & The likelihood assumes Gaussian innovations. \\
Deterministic trends in the measurement equation & Not handled by the
inversion. \\
Diffuse and heteroskedastic filters & Not applicable; $\Sigma_\varepsilon$ is
assumed constant. \\
\option{forecast} & The smoother does not populate the \field{oo\_.Smoother}
fields \texttt{forecasts.run} requires. \\
\command{histval} & The simulation always starts from the steady state. \\
\bottomrule
\end{longtable}

\noindent
A second group produces a warning rather than an error, because the run is
valid but an output is silently missing or not comparable:
\option{filtered\_vars}, \option{filter\_covariance},
\option{filter\_decomposition}, \option{moments\_varendo} and
\option{bayesian\_irf} rely on Kalman-filter objects the inversion does not
produce; \option{hp\_filter} and \option{bandpass} are honoured by the RE
column of the moment tables but not by the simulated SL column.

\paragraph{\option{prefilter}} is supported: the measurement constant is
dropped when the data have been demeaned.

\paragraph{Octave.} The learning streams require \matlab{}
\texttt{RandStream} substreams, which have no Octave equivalent, and both
\texttt{simult\_SL} and \texttt{IF\_SL} reject a non-\texttt{RandStream}
argument. The package therefore does not run under Octave at all.


%================================================================
%  SECTION 6 — STORED RESULTS
%================================================================
\section{Stored results reference}
\label{sec:output}

\subsection{\texttt{oo\_} fields}

\begin{longtable}{>{\ttfamily\small}p{6.2cm} p{8cm}}
\toprule
\normalfont\textbf{Field} & \textbf{Description} \\
\midrule
\endhead
oo\_.dr.ghr & $R$ matrix, $n_{\text{endo}} \times n_{\text{endo}}$ \\
\midrule
\multicolumn{2}{l}{\normalfont\textit{Simulation}} \\
\midrule
oo\_.endo\_simul & Simulated path, $n_{\text{endo}} \times T$ \\
oo\_.learning.social.beliefs & Post-tournament beliefs, $n_{\text{fwrd}} \times J \times T$ \\
oo\_.learning.social.beliefs\_candidate & Post-injection candidates, same size \\
oo\_.learning.social.fitness & Fitness values, same size \\
oo\_.learning.social.expectations & Aggregate wedge $\varphi_t + O_r\varepsilon_t$, $n_{\text{fwrd}} \times T$ \\
oo\_.learning.social.expectations\_ahead & $h$-step-ahead expectations, $n_{\text{fwrd}} \times T \times$ \texttt{t\_max} \\
oo\_.learning.social.final\_beliefs & Beliefs at $t = T$, $n_{\text{fwrd}} \times J$ \\
oo\_.learning.social.y\_RE & RE path under the same structural shocks \\
oo\_.learning.social.decomp & Historical decomposition, $n_{\text{endo}} \times T \times n_{\text{chan}}$ \\
\midrule
\multicolumn{2}{l}{\normalfont\textit{Moments and IRFs}} \\
\midrule
oo\_.mean, oo\_.var, oo\_.autocorr & Monte Carlo moments \\
oo\_.learning.social.moments.mean\_ci & Empirical intervals on the means \\
oo\_.learning.social.moments.std\_ci & Empirical intervals on the standard deviations \\
oo\_.learning.social.moments.cross\_section & Cross-sectional belief statistics \\
oo\_.learning.social.moments.mc\_raw & Replication-level raw data \\
oo\_.learning.social.irf.\textit{shock} & \texttt{.median}, \texttt{.mean}, \texttt{.bands} \\
\midrule
\multicolumn{2}{l}{\normalfont\textit{Estimation}} \\
\midrule
oo\_.SmoothedShocks & Smoothed shocks, belief channels included \\
oo\_.SmoothedVariables & Smoothed states \\
oo\_.Smoother.Constant & Constant removed from the measurement equation \\
oo\_.learning.social.smoother & \texttt{.expectations}, \texttt{.forward\_states}, \texttt{.forward\_vars}, \texttt{.news} \\
\midrule
\multicolumn{2}{l}{\normalfont\textit{Adaptive learning}} \\
\midrule
oo\_.learning.adaptive.info & \texttt{.beta\_hist}, \texttt{.beta\_final}, \texttt{.pf\_hits} \\
oo\_.learning.adaptive.moments & Same structure as the SL moments \\
\bottomrule
\end{longtable}


\subsection{\texttt{options\_.learning.social} fields}

\begin{longtable}{>{\ttfamily\small}p{4.4cm} >{\ttfamily\small}p{3cm} p{6.6cm}}
\toprule
\normalfont\textbf{Field} & \normalfont\textbf{Type} & \textbf{Description} \\
\midrule
\endhead
N\_agents  & integer                  & Population size $J$ \\
mut\_p     & double [nfwrd$\times$1]  & News probability $\mu$ \\
wed\_p     & double                   & Tournament participation share \\
mut\_sd    & double [nfwrd$\times$1]  & Belief dispersion $\xi$ \\
rhoGN      & double [nfwrd$\times$1]  & Fitness memory $\delta$ \\
t\_max     & integer                  & Horizon of the reported expectations \\
Or         & double [nfwrd$\times$nexo] & Belief channel selection \\
Ox         & double [nexo$\times$nexo]  & Structural shock projector $I - O_r'O_r$ \\
id\_SL     & integer vector           & Indices of the learned variables \\
n\_SL      & integer                  & Number of learned variables \\
decomp     & logical                  & Compute the historical decomposition \\
irf\_quantiles & double vector        & GIRF band levels \\
\midrule
\multicolumn{3}{l}{\normalfont\textit{Added by \texttt{setup\_learning\_estimation}}} \\
\midrule
O          & double [nobs$\times$nendo] & Observation selection \\
id\_obs    & integer vector           & Indices of the observed variables \\
id\_fwrd   & integer vector           & Indices of the forward variables \\
nfwrd      & integer                  & Number of forward variables \\
omega      & double [nfwrd$\times$nendo] & Forward selection \\
mut\_pid, mut\_vid & integer vectors  & Parameter and variable index maps for $\xi$ \\
rhoGN\_pid, rhoGN\_vid & integer vectors & Idem for $\delta$ \\
mutp\_pid, mutp\_vid & integer vectors & Idem for the per-variable $\mu$ \\
estimation & logical                  & Set when the estimation objects are built \\
\bottomrule
\end{longtable}


%================================================================
%  APPENDIX
%================================================================
\appendix

\section{Internal architecture}
\label{app:architecture}

\subsection{Package map}
\label{app:architecture:map}

\file{+learning} holds what is shared between mechanisms, and one sub-package
per mechanism.

\begin{longtable}{>{\ttfamily\small}p{5.5cm} p{8.5cm}}
\toprule
\normalfont\textbf{File} & \textbf{Role} \\
\midrule
\endhead
\multicolumn{2}{l}{\normalfont\textit{+learning}} \\
\midrule
learning\_streams.m & Build the belief and shock substreams from one seed \\
is\_active.m        & True when a learning mechanism is configured \\
getopt.m            & Read an optional field of an option structure, with a default. It resolves fields via \texttt{isfield} and must never be used to read a \file{.mod} parameter: a parameter is an entry of \field{M\_.param\_names}, never a field of \field{M\_}, so \texttt{learning.getopt(M\_, \ldots)} silently returns the default \\
\midrule
\multicolumn{2}{l}{\normalfont\textit{+learning/+social --- simulation}} \\
\midrule
setup\_options\_SL.m & Build \field{options\_.learning.social} from the \texttt{SL\_*} parameters \\
check\_support.m     & Reject or flag unsupported configurations \\
warn\_wedge\_variance.m & Report \texttt{a\_VAR} shocks left with a non-zero variance while the mechanism is off; called from the dispatch, not from the package \\
compute\_ghr.m       & Compute and validate the $R$ matrix \\
simult\_SL.m         & Forward simulation with mutation and tournament \\
moments\_SL.m        & Monte Carlo moments with empirical intervals \\
disp\_moments\_SL.m  & Moment tables \\
irf\_SL.m            & Paired Monte Carlo generalised IRFs \\
plot\_decomp.m       & Historical decomposition plots \\
\midrule
\multicolumn{2}{l}{\normalfont\textit{+learning/+social --- estimation}} \\
\midrule
setup\_learning\_estimation.m & Build the $\theta$-invariant estimation objects, run once \\
refresh\_learning\_params.m   & Inject the current draw, hot path \\
dsge\_likelihood\_SL.m        & Log-posterior, native signature \\
IF\_SL.m                      & Inversion filter core \\
store\_SL\_smoother.m         & Route the smoother payload into \field{oo\_} \\
plot\_SL\_smoother\_results.m & Post-estimation figures \\
\midrule
\multicolumn{2}{l}{\normalfont\textit{+learning/+adaptive}} \\
\midrule
setup\_options\_AL.m  & Build \field{options\_.learning.adaptive} \\
check\_support\_AL.m  & Rejected configurations \\
init\_beliefs\_AL.m   & Model-consistent initial beliefs \\
update\_beliefs\_AL.m & Constant-gain RLS updating \\
resolve\_alm\_AL.m    & Map beliefs into the actual law of motion \\
simult\_AL.m          & Forward simulation with belief updating \\
moments\_AL.m         & Monte Carlo moments \\
disp\_moments\_AL.m   & Moment tables \\
\bottomrule
\end{longtable}


\subsection{Entry-point convention}

Every mechanism entry point takes \texttt{(M\_, options\_, oo\_, \ldots)} and
returns \field{oo\_}:

\begin{lstlisting}[style=matlab]
oo_ = learning.social.simult_SL(M_, options_, oo_, ee, belief_rng);
oo_ = learning.adaptive.simult_AL(M_, options_, oo_, ee, belief_rng);
oo_ = learning.social.moments_SL(M_, options_, oo_);
oo_ = learning.adaptive.moments_AL(M_, options_, oo_);
\end{lstlisting}

Shock matrices are \texttt{[T x nexo]} in both packages. \texttt{simult\_AL}
accepts a stream argument it does not use, so that the two mechanisms share one
signature.


\subsection{Hook points in Dynare core}
\label{app:architecture:hooks}

Two core files are modified.

\paragraph{\file{stoch\_simul.m}} Five points: the option scan at the top,
\texttt{compute\_ghr} after \texttt{resol}, the simulation block, the moments
block and the IRF loop. The shock draw is shared between mechanisms; only the
final dispatch differs:

\begin{lstlisting}[style=matlab]
if learning.is_active(options_)
    id_struct   = find(diag(M_.Sigma_e) > 0);
    cs          = get_lower_cholesky_covariance(M_.Sigma_e, ...
                      options_.add_tiny_number_to_cholesky);
    chol_struct = cs(id_struct, id_struct)';
    [belief_rng, shock_rng] = learning.learning_streams( ...
        options_.DynareRandomStreams.seed, ...
        options_.DynareRandomStreams.algo, 'sim', 1);
    ee = zeros(options_.periods, M_.exo_nbr);
    ee(:, id_struct) = randn(shock_rng, options_.periods, ...
                             numel(id_struct)) * chol_struct;
    if learning.is_active(options_, 'social')
        oo_ = learning.social.simult_SL(M_, options_, oo_, ee, belief_rng);
    else
        oo_ = learning.adaptive.simult_AL(M_, options_, oo_, ee, belief_rng);
    end
end
\end{lstlisting}

\noindent
The option scan at the top of the file is the activation dispatch. It reads the
two master switches from \field{M\_.params}.

\begin{lstlisting}[style=matlab]
idx = find(strcmp(M_.param_names, 'SL_active'), 1);
if ~isempty(idx) && isfinite(M_.params(idx)) && M_.params(idx) ~= 0
    options_ = learning.social.setup_options_SL(M_, options_);
else
    learning.social.warn_wedge_variance(M_);
end
idx = find(strcmp(M_.param_names, 'AL_active'), 1);
if ~isempty(idx) && isfinite(M_.params(idx)) && M_.params(idx) ~= 0
    options_ = learning.adaptive.setup_options_AL(M_, options_);
end
\end{lstlisting}



\paragraph{\file{dynare\_estimation\_1.m}} Three points: the activation
dispatch, the choice of the objective function and the smoother branch.

\begin{lstlisting}[style=matlab]
% activation dispatch, mirroring stoch_simul
idx = find(strcmp(M_.param_names, 'SL_active'), 1);
if ~isempty(idx) && isfinite(M_.params(idx)) && M_.params(idx) ~= 0
    options_ = learning.social.setup_learning_estimation(M_, options_, estim_params_);
end
\end{lstlisting}

\begin{lstlisting}[style=matlab]
% objective function
elseif learning.is_active(options_, 'social') && ...
        learning.getopt(options_.learning.social, 'estimation', false)
    objective_function = str2func('learning.social.dsge_likelihood_SL');
\end{lstlisting}

\medskip
\noindent
Dynare reaches the classical smoother at a different point in the code:

\begin{lstlisting}[style=matlab]
if learning.is_active(options_, 'social') && ...
        (learning.getopt(options_, 'smoother', false) || ...
         learning.getopt(options_, 'frequentist_smoother', false))
    [~,~,~,~,~,model_solution] = learning.social.dsge_likelihood_SL(...);
    oo_ = learning.social.store_SL_smoother(oo_, M_, options_, model_solution);
    learning.social.plot_SL_smoother_results(oo_, M_, options_, dataset_);
elseif options_.frequentist_smoother
    oo_ = save_display_classical_smoother_results(...);
end
\end{lstlisting}


\subsection{Test suites}
\label{app:architecture:tests}

\file{tests/learning/social} holds seven \file{.mod} files, each Dynare file coresponds to a serie of tests, running in about eight seconds in total.
\file{tests/learning/adaptive} holds one. Both use a shared file in
\file{tests/learning} and are documented in a manifest per mechanism, which
lists what each check asserts, the reference values, and what is missing. 

The assertions that matter most:

\begin{itemize}[leftmargin=2em]
\item \textbf{Simulation--inversion round trip.} \texttt{IF\_SL} recovers the
      shocks, the state path and the belief path produced by
      \texttt{simult\_SL}, to $10^{-12}$.
\item \textbf{RE limit.} With zero belief dispersion and no news, the
      belief wedge vanishes and \texttt{simult\_SL} reproduces the native
      \texttt{simult\_} exactly; the generalised IRFs reproduce the analytic
      responses.
\item \textbf{Deterministic likelihood.} Two evaluations at the same $\theta$
      return bit-identical values.
\item \textbf{Steady-state invariance.} Beliefs are invariant to a shift of the
      steady state while the path shifts by exactly that amount.
\item \textbf{Smoother dispatch.} After an estimation,
      \field{oo\_.learning.social.smoother} exists and
      \field{oo\_.UpdatedVariables}, which marks the native Kalman path, does
      not.
\item \textbf{Input guards.} Thirty-two rejected configurations on the social
      side and nine on the adaptive one, each checked on
      \texttt{ME.identifier} 
\item \textbf{Selection follows the flags.} With a single
      \texttt{SL\_learn\_VAR} left at 1, exactly one variable is learned.
\end{itemize}

\end{document}